\section{Conclusion and Future Work}
\label{chap:conclusion}

In this paper, we have described the first mechanization of mailbox types and
their semantics, as well as the syntax and semantics of the Pat programming
language, in the Rocq Prover.
%
We show that mechanizing Pat is nontrivial: in particular, environment
combinations and subtyping pose a big challenge due to the numerous
technical lemmas needed.
%
Our work successfully provides a significant number of
machine-checked proofs for various properties about mailbox types and Pat's
static semantics.  During the mechanization process, we also identified and
fixed multiple oversights in the original definitions and proofs.

The next step for future work is the mechanization of the operational semantics,
which would allow us to formulate machine-checked proofs of properties such as
preservation and progress.
%
Initial work has shown that a primary difficulty in formalizing Pat's
operational semantics is the constant manipulation of de Bruijn indices.  While
tedious, previous work \cite{zalakainLeftoversMechanisationAgda2021} was able to
successfully formalize a similar semantics with this representation of
variables.
%
Alternatively, we could make heavier use of Rocq's dependent types.  In some
mechanizations of the $\pi$-calculus \cite{ambalHOPiCoq2021}, the type system
keeps track of available channel endpoints for each process, which could
potentially reduce the proofs regarding de Bruijn indices.

% A further research area, after successfully proving preservation and progress, is the verification of the algorithmic type system of Pat.
% This thesis focuses on the declarative type system as a base of correctness.
% However, Pat's implementation makes use of the algorithmic type system.
% Formally verifying whether both type systems provide the same properties would be an interesting further work.
